\section{Pruning and unique recovery of long paths}
\label{prune:section}

To pass from the two towers to one graph, we need both a small deletion
fraction in every layer and long words that determine their paths.
We first prove that a word recovers padded stream prefixes at interior
positions in \(Y^\#\).  This does not yet identify the layer: the labels
determine its sign and residue but need not determine its absolute value,
and padding can conceal a shorter stream.  The subsequent deletions
resolve this ambiguity, make the gated and auxiliary graphs agree, and
remove short cyclic paths.

A path in this section is a finite edge walk,
possibly with repeated vertices.  Its radius at a specified position
is the smaller of the numbers of edges available in the two directions
from that position.  A nonbacktracking path never follows an edge
immediately by its reverse.  A closed path is cyclically
nonbacktracking when this also holds at its cyclic junction.

\begin{lemma}[Recovery of stream prefixes]
\label{prune:prefix-recovery}
Suppose two nonbacktracking paths in \(Y^\#\) have the same oriented
label word.  At corresponding \(A\)-positions of radius at least
\(2j+1\), both padded streams agree through length \(r+sj\), for
every \(j\geq0\).  At corresponding \(B\)-positions of radius at
least \(2j\), both padded streams agree through length \(sj\), for
every \(j\geq1\).  The compared positions need not be in the same
layer.
\end{lemma}

\begin{proof}
The orientation of a letter identifies the side of either endpoint.
Reading its underlying positive label gives the family, the serial
number and \(x=b\Vert x'\), and therefore both slots.  In particular,
the same word identifies the same slot names, diagonal coefficients,
spill maps, and family rules at corresponding positions.  It also
identifies the signs and residues modulo \(K\) of the layer indices,
although their absolute values may differ.

At an interior position the two incident slots are distinct.  Otherwise
the incoming and outgoing edges would be the unique edge at the same
slot, so the path would backtrack.  Write their evaluations on a current
block as
\[
 Q^a_h=v_h+\lambda_a u_h+S_a u_{h-1},\qquad
 Q^b_h=v_h+\lambda_b u_h+S_b u_{h-1}.
\]
Here \(u_h,v_h\) are the full padded current blocks, multiplication
uses their corresponding field, and the last terms are the spills
from the preceding block, extended by zero as needed.  For the first
block there is no preceding spill.  If the two evaluation blocks and
the preceding input block have been recovered, subtracting gives
\begin{equation}
 (\lambda_a-\lambda_b)u_h
   =Q^a_h-Q^b_h-(S_a-S_b)u_{h-1}.
 \label{prune:block-inversion}
\end{equation}
The coefficient difference is nonzero and hence invertible.  Thus
\(u_h\), followed by \(v_h\), is recovered.  Different spill
lengths for the two incident slots do not affect this argument.

Call the asserted \(A\)-statement at index \(j\) the assertion
\(\mathcal A_j\), and the \(B\)-statement \(\mathcal B_j\).
Interpret \(\mathcal B_0\) as the vacuous recovery of the empty
prefix.  At radius one on side \(A\), the first \(r\) bits of each
incident evaluation are present in its label \(x\), because
\(p=r+\alpha\geq r\).  Equation~\eqref{prune:block-inversion}
for the first block proves \(\mathcal A_0\).

Suppose \(j\geq1\).  At a \(B\)-position of radius \(2j\), both
neighbors satisfy \(\mathcal A_{j-1}\); centrally we already have
\(\mathcal B_{j-1}\).  For either incident edge, the second equation
of \eqref{tow:edge-equations} shows that to know \(Q_B\) through
length \(sj\), beyond the label prefix \(x'\), it suffices to know
the \(\mathcal J\)-output through length
\[
 sj-p'=r+s(j-1)-\beta\leq r+s(j-1).
\]
All source \(A\)-blocks involved are known in full at the neighbor.
Every key for the transformation of \(A\)-block \(j-1\), or an
earlier block, uses central \(B\)-blocks at most \(j-1\), together
with earlier blocks of that neighbor's \(A\)-stream.  These are all
known.  Consequently the indicated output prefix, including every
gate decision, agrees in the two paths.  Both central evaluation
prefixes therefore agree through \(sj\).  Their preceding spill
inputs are already known by \(\mathcal B_{j-1}\), so
\eqref{prune:block-inversion} recovers the next full \(B\)-block.
This proves \(\mathcal B_j\).

At an \(A\)-position of radius \(2j+1\), both neighbors satisfy
\(\mathcal B_j\), and centrally we have
\(\mathcal A_{j-1}\).  By the first equation of
\eqref{tow:edge-equations}, the evaluation prefix through \(r+sj\)
needs only the \(\mathcal F\)-output through
\[
 r+sj-p=sj-\alpha\leq sj.
\]
The required neighbor \(B\)-blocks through block \(j\) are fully
known.  The keys for such a block use central \(A\)-blocks only
through block \(j-1\), and preceding neighbor \(B\)-blocks.
These again are known, so both evaluation prefixes agree.  Subtract
the already known preceding spill and apply
\eqref{prune:block-inversion}.  This proves
\(\mathcal A_j\) and completes the alternating induction.

This proof uses full padded blocks.  A permutation may mix all the
bits of a block, but its entire source block is already known at the
neighbor before it is used.  Inactive bits cause no difficulty because
the equations are padded-stream equations.  When a rule lies beyond
one of the compared paths' actual cutoffs, it is still the same
family rule on both paths; equality of the preceding input blocks
gives equality of its key and gate.  This is precisely why
Lemma~\ref{prune:prefix-recovery} is proved in \(Y^\#\), with its
family-wide gated rules.  The
inequalities used above also include \(\alpha=0\), \(\beta=0\),
and a final partial block of any length \(1\leq t\leq s\).
\end{proof}

\begin{corollary}[Few short cyclic paths at one layer]
\label{prune:cycle-count}
Fix a layer and a side.  A nonempty cyclically reduced oriented label
word is the word of at most one based cyclically nonbacktracking
closed path with its base on that layer and side.  Consequently the
number of vertices of \(A_i\), or of \(B_i\), supporting such a
closed path of length at most \(g_i\) is at most
\[
 \sum_{a=1}^{\lfloor g_i\rfloor}(2m_\ell)^a.
\]
For all sufficiently large \(\ell\), this is a fraction at most
\(2^{-0.9H_i}\) of either vertex set, uniformly in \(i\).
\end{corollary}

\begin{proof}
Repeat two closed paths with the same word arbitrarily many times
in both directions.  Cyclic nonbacktracking ensures that every such
repetition is nonbacktracking.  At corresponding base positions,
Lemma~\ref{prune:prefix-recovery} recovers arbitrarily long prefixes.
Since the layer and side are fixed, this eventually recovers all
active coordinates of both vertices, making their base states equal.
Foldedness then determines their entire based paths.  Conversely a
cyclically nonbacktracking closed path has a cyclically reduced word:
inverse successive letters, including at the cyclic junction, would
force a backtrack by foldedness.

There are at most \((2m_\ell)^a\) words of length \(a\); counting
all words only enlarges the bound.  Both vertex-set sizes are at least
\(2^{H_i}\), and \eqref{tow:alphabet} gives
\[
 \frac{\log_2(2m_\ell)}{100\log_2 d}<\frac1{20}
\]
for large \(\ell\).  If \(\lfloor g_i\rfloor\geq1\), the
geometric sum is at most \(2(2m_\ell)^{g_i}\), so its fraction is
at most \(2^{1-0.95H_i}\leq2^{-0.9H_i}\).  If the sum is empty
the claim is immediate.  Since \(H_i\geq H_0\to\infty\), the
threshold is uniform over all layers.
\end{proof}

We now specify the deleted sets.  They are sets of vertices of the
common vertex set of \(Y^0\) and \(Y^\#\).

\paragraph{Gate deletions.}
At each vertex and each adjacent-pair slot evaluate the gate at that
pair's own final block, whether or not the current auxiliary rule
actually consults it.  Delete the vertex if any of these gates is
closed.  The inputs before the last block are determined without
using its current permutation, so this is a well-defined test on the
slot state, identical for the gated and auxiliary prescriptions.

At a fixed \(A\)-slot, conditional on \(x\), the coordinates
\(u_A\) and the tail of \(Q_A\) are independent uniform bits,
by the parametrization in Lemma~\ref{tow:matching-clusters}.
The triangular equations turn them bijectively into the interleaved
input \(u\)-values.  Thus all preceding input coordinates are
independent uniform bits.  The same reasoning at a fixed \(B\)-slot
uses \(x'\), its fixed bucket, \(u_B\), and the tail of \(Q_B\).
In particular the four full \(s\)-blocks preceding the cutoff are
jointly zero with probability exactly \(2^{-4s}\) at either slot.
The union bound and the maximum degree \(6D\) show that these
deletions have fraction at most
\begin{equation}
 \varepsilon_{\rm gate}=6D\,2^{-4s}.
 \label{prune:gate-budget}
\end{equation}

On vertices surviving these deletions the induced graphs from
\(Y^0\) and \(Y^\#\) are identical.  Indeed, from a surviving
slot state its current final gate is open, so the two endpoint
constructions agree.  All lower rules already agree.  This argument
works from either side of an edge and proves equality of the induced
edge sets, including multiplicities and labels.

\paragraph{Separation of heights.}
The next deletion distinguishes larger-layer vertices from the zero
padding of smaller-layer vertices with the same sign and residue.
For a vertex on either side at layer \(i\), consider every layer
\(i'\) with
\[
 \operatorname{sgn}(i')=\operatorname{sgn}(i),\qquad
 |i'|\equiv|i|\pmod K,\qquad |i'|<|i|.
\]
Put \(m=(|i|-|i'|)/K\geq1\) and
\begin{equation}
 q(i,i')=\min\{H_{i'},mK/2\}.
 \label{prune:zero-window}
\end{equation}
Delete the vertex if all its own \(u\)-bits immediately after that
side's cutoff at \(i'\), for these \(q(i,i')\) positions, are zero.
This interval consists of active bits at \(i\), because on either
side the length increases by exactly \(mK/2\) between these two
layers.  Its zero probability under uniform vertex measure is
\(2^{-q(i,i')}\).  No layers of distinct signs are tested, and
zero has no smaller layer with its sign.

The union bound gives, on every side and layer,
\begin{equation}
\begin{split}
 \varepsilon_{\rm sep}
 &\leq\sum_{m\geq1}2^{-mK/2}+\sum_{i'\in\mathbb Z}2^{-H_{i'}}\\
 &=\frac{2^{-K/2}}{1-2^{-K/2}}+3\,2^{-H_0}.
\end{split}
\label{prune:separation-budget}
\end{equation}
Here \(2^{-\min(a,b)}\leq2^{-a}+2^{-b}\), and the second sum
is evaluated using \(H_{i'}=H_0+|i'|\).  Thus the estimate stays
uniform even when the larger layer has arbitrarily many smaller
layers of its residue.

\paragraph{Short-cycle deletions.}
At layer \(i\), delete every vertex supporting a nonempty cyclically
nonbacktracking closed path in \(Y^\#\) of length at most \(g_i\).
By Corollary~\ref{prune:cycle-count}, the fraction deleted on either
side is at most \(2^{-0.9H_i}\).  This test is performed in the
already specified graph \(Y^\#\), so it does not depend circularly
on the final induced graph.

Let \(Y\) be the graph induced on the vertices surviving all three
deletions.  It is an induced subgraph both of \(Y^0\) and of
\(Y^\#\).  Combining the three budgets gives a uniform bound
\begin{equation}
\begin{split}
 \varepsilon_\ell
 &=6D\,2^{-4s}
   +\frac{2^{-K/2}}{1-2^{-K/2}}
   +3\,2^{-H_0}+2^{-0.9H_0}\longrightarrow0.
\end{split}
\label{prune:total-budget}
\end{equation}
For example, \(D2^{-4s}=2^{-32\ell}\).  Fractions here and below
are always taken relative to the original sets \(A_i,B_i\), not
relative to an already pruned set.

\begin{proposition}[The tower properties]
\label{prune:tower-properties}
For all sufficiently large \(\ell\), the preceding choices give a
graph \(Y^0\) on sets of sizes
\(|A_i|=d2^{H_i}\) and \(|B_i|=2^{H_i}\), and its induced
subgraph \(Y\), with these properties.
\begin{enumerate}
\item For every \(|i-k|\leq1\), the \(Y^0\)-block between
\(A_i\) and \(B_k\) has degrees \(d,D_{ik}\) with
\(D_{ik}=D2^{|i|-|k|}\).  Its counting incidence divided by
\(\sqrt{dD_{ik}}\) maps normalized constants to normalized
constants, as does its adjoint.  On their orthogonal complements its
norm is at most \(1-\chi\), with the absolute constant \(\chi\)
of Proposition~\ref{exp:realized-gap}.
\item The fraction removed from any \(A_i\) or \(B_i\) is at
most \(\varepsilon_\ell\) of \eqref{prune:total-budget}.
\item The labels of \(Y\) are folded, with positive alphabet size
at most \(m_\ell=O(d2^r(1+s)^2)\).
\item Every nonbacktracking path in \(Y\) beginning at layer \(i\)
and having length at least \(L_i\) is determined by its oriented
word among all paths in \(Y\), including paths beginning in other
layers.
\item There is no nonempty cyclically nonbacktracking closed path
in \(Y\), based at layer \(i\), of length at most \(g_i\).
\end{enumerate}
\end{proposition}

\begin{proof}
Only the fourth item requires further proof.  The degree and norm
claims are established in \eqref{tow:degrees} and
Proposition~\ref{exp:realized-gap}.  For normalized counting
constants, for example, the constant entry after applying normalized
incidence is
\(d/(\sqrt{dD_{ik}}\sqrt{|B_k|})=1/\sqrt{|A_i|}\),
using \(d|A_i|=D_{ik}|B_k|\); the adjoint calculation is the same.
The deletion and alphabet claims have already been proved.
Any short cyclic path in \(Y\) is also one in \(Y^\#\), and its
base vertex was deleted in the third step, proving the fifth item.

Take a path beginning in layer \(i\), and compare it with a path
having the same word.  By foldedness, any path with this word is also
nonbacktracking.  It suffices to identify the segment of length
\(N=\lceil L_i\rceil\) at the beginning of each path.  Choose a
midpoint position, whose radius is at least \(\lfloor N/2\rfloor\),
and let its layer on the first path be \(j\).  Every edge changes
the layer by at most one, so
\[
 H_j\leq H_i+N.
\]
Lemma~\ref{prune:prefix-recovery} recovers both padded streams at
that position through length at least
\[
 \frac{sN}{4}-2s
 \geq\frac{25s}{\log_2d}H_i-2s
 =62.5H_i-2s.
\]
Since \(H_i\geq H_0\) and
\(N\leq100H_i/\log_2d+1\), this quantity is greater than
\(2H_j\) for all sufficiently large \(\ell\), uniformly in \(i\).
For side \(A\) the exact recovered length at radius \(R\) is
\(r+s\lfloor(R-1)/2\rfloor\); on side \(B\) it is
\(s\lfloor R/2\rfloor\).  These formulas give the displayed
lower bound on both sides, including either parity of \(N\).

The common incident labels force the two midpoint layers to have the
same sign and residue modulo \(K\).  If the layer indices are equal,
all active coordinates are recovered, since each has length at most
\(H_j\), and the midpoint vertices coincide.
Suppose instead they differ.  Write \(i_-\) and \(i_+\) for their
indices ordered by absolute value, and put \(H_-=H_{i_-}\).
Then \(H_-\leq H_j\) and \(|i_+|-|i_-|=mK\) for some
\(m\geq1\).  If \(l_-\) is the smaller vertex's own
\(u\)-length, the interval tested at the larger vertex in
\eqref{prune:zero-window} ends no later than
\[
 l_-+\min\{H_-,mK/2\}\leq 2H_-\leq2H_j.
\]
The smaller vertex's padding is zero on this entire interval.
Prefix recovery forces the larger vertex's active \(u\)-bits on
the same interval to be zero.  It would therefore have been removed
by the second deletion, a contradiction.  This reasoning also covers
the case where the first path has the larger midpoint height.  The
case of a zero midpoint index was already in the equal-layer case,
since its sign is distinguished.

Thus the midpoint vertices coincide in all cases.  Foldedness
determines an edge uniquely from its initial vertex and its label;
applying this forward and backward from the common midpoint identifies
the whole two segments.  Continuing forward identifies any remaining
part of the compared words.  This proves the fourth item.
\end{proof}

\begin{corollary}[Effect of deletion on constant spaces]
\label{prune:constant-compression}
Let \(\Pi\) be the orthogonal projection that retains the vertices
of \(Y\) in the counting Hilbert space of the full tower, with any
additional Hilbert-space tensor factor.  Let \(C\) be the subspace
of vectors that are constant on each separate \(A_i\) and \(B_i\)
in the vertex coordinate.  Then
\[
 \|(1-\Pi)|_C\|\leq\sqrt{\varepsilon_\ell}.
\]
In particular the same bound holds on every subspace of \(C\),
uniformly in the additional tensor factor.
\end{corollary}

\begin{proof}
On one normalized constant vector the squared mass removed is exactly
the fraction of vertices removed from that set.  Distinct side and
layer sets are orthogonal, so the squared losses add, and each is
bounded by \(\varepsilon_\ell\) times the corresponding squared
input norm.  Tensoring changes neither calculation.
\end{proof}
